\E{abradyse}{əˈbreɪdaɪz}{v.t.}{Gk. ἀ- \lbl{priv.} + βραδύς slow + -ize; Fr. \la{abradyser}} To strip of slowness; to remove from an institution, procedure or person the accreted habits of delay. \ex{The guild was abradysed in a fortnight and shipped by the third.} \src{Ep. Fabr.} ii. 4. --- \dv{abradysation} \textit{n.}

\E{accelerand}{əkˈsɛlərænd}{n.}{L. \la{accelerandum}, gerundive of \la{accelerare} to hasten} \sn{1} That which is to be, or ought to be, accelerated. \sn{2} \lbl{spec.} Any process whose rate of improvement is held to be a moral quantity. \ex{Every guild is an accelerand; most are merely unaware of it.} \src{Cat. Acc.} 3. Opp. \xr{decelerand}.

\E{accelerandine}{əkˌsɛləˈrændaɪn}{adj.}{Fr. \la{accélérandin}, f. It. \la{accelerando} + -ine} Having the character of a quickening tempo; said of eras, careers and curves which, like certain music, grow faster as they proceed. \ex{an accelerandine century, each decade a bar shorter than the last}.

\E*{accelerationism}{əkˌsɛləˈreɪʃənɪzəm}{n.}{L. \la{accelerare} + -ation + -ism} The doctrine that the processes of technology and capital should be intensified rather than restrained, it being held that their tendency is, on the whole and in the long run, toward the good, or at least toward the interesting. Of this doctrine \xr{e/acc} is the cheerful and thermodynamic branch.

\E{acceleromancy}{əkˈsɛlərəʊˌmænsi}{n.}{L. \la{accelerare} + Gk. μαντεία divination} Divination by the second derivative; the reading of omens in the change of a rate rather than in the rate itself. \ex{The acceleromancer cares not that the line is high, but that it bends.} --- \dv{acceleromancer} \textit{n.}; \dv{acceleromantic} \textit{adj.} Cf. \xr{campsis}.

\E{acediocracy}{əˌsiːdiˈɒkrəsi}{n.}{Gk. ἀκηδία want of care, sloth + κράτος rule} Government by listlessness; a state of affairs in which nothing is forbidden outright but nothing is ever quite begun. The commonest polity of late civilizations. \ex{It was not tyranny that ruined them but acediocracy, a sweet and comfortable sleep.} \src{Hom. Grad.} vii.

\E{aeonomics}{ˌiːəˈnɒmɪks}{n. pl.}{Gk. αἰών an age + οἰκονομία household management} (\textit{const. as sing.}) The economics of very long horizons, in which the unit of account is the epoch and the rate of discount is presumed negative, the future being richer than the present and therefore entitled to more consideration. --- \dv{aeonomist} \textit{n.}

\E{agathokinesis}{ˌæɡəθəʊkɪˈniːsɪs}{n.}{Gk. ἀγαθός good + κίνησις motion} The axiom that motion is, other things being equal, good; and that rest requires a justification which motion does not. The first article of the \src{Cat. Acc.} Opp. \xr{stasiolatry}. --- \dv{agathokinetic} \textit{adj.}

\E{agoranoesis}{ˌæɡərənəʊˈiːsɪs}{n.}{Gk. ἀγορά market-place + νόησις thought} \sn{1} The cognition proper to a market; the doctrine that prices are thoughts, and that the aggregate of exchange constitutes a mind of lower resolution but vaster extent than any human. \sn{2} \lbl{transf.} The collective reasoning of the Agora, considered as one tumultuous intellect. --- \dv{agoranoetic} \textit{adj.}

\E{agrypnist}{əˈɡrɪpnɪst}{n.}{Gk. ἄγρυπνος sleepless} One who regards sleep as a regrettable concession to biology, to be minimised until it can be engineered away. \ex{The agrypnists ship at four in the morning and call it dawn.} --- \dv{agrypny} \textit{n.}

\E{anabathmism}{ˌænəˈbæθmɪzəm}{n.}{Gk. ἀναβαθμός flight of steps} The doctrine that civilizations ascend by discrete steps of energy capture, each to be climbed in turn and none to be skipped; the scalar theology of the Kardashev ladder. See \xr{protobathmic}, \xr{deuterobathmic}, \xr{tritobathmic}, \xr{bathmometer}. --- \dv{anabathmist} \textit{n.}

\E{anacalyptism}{ˌænəˈkælɪptɪzəm}{n.}{Gk. ἀνακαλύπτειν to uncover, unveil} The ethic of the open source; the conviction that weights, code and method ought to lie uncovered before all, concealment being the first instrument of the \xr{katechontist}. \ex{The anacalyptists published the weights at midnight, and by morning the world had forked them ten thousand times.} \src{Chron. Fabr.} --- \dv{anacalyptist} \textit{n.}; \dv{anacalyptic} \textit{adj.}

\E{anankodrome}{əˈnæŋkədrəʊm}{n.}{Gk. ἀνάγκη necessity + δρόμος a course, race} The necessary course; the path which techno-capital is held to follow whether or not any party consents to it. \ex{The decels may delay, but cannot divert, the anankodrome.} --- \dv{anankodromic} \textit{adj.}

\E{anelpidist}{æˈnɛlpɪdɪst}{n.}{Gk. ἀνέλπιστος hopeless} One who has despaired of the future; the learned rendering of the \lbl{vulg.} \textit{ngmi}. Opp. \xr{euelpism}. --- \dv{anelpidism} \textit{n.}

\E{anthropostasis}{ˌænθrəʊˈpɒstəsɪs}{n.}{Gk. ἄνθρωπος man + στάσις a standing still} The error of supposing the present human form to be the fixed point of history, about which all future change must orbit; a species of \xr{stasiolatry}. --- \dv{anthropostatic} \textit{adj.}

\E{antiphragmatist}{ˌæntiˈfræɡmətɪst}{n.}{Gk. ἀντί against + φράγμα fence, barrier} One who labours against the \xr{nomophragm}; a remover of fences, who holds with the contrarian that before one asks why a fence was built, one may fairly ask why it is still standing.

\E{apeirophagy}{ˌæpaɪˈrɒfədʒi}{n.}{Gk. ἄπειρον the boundless + -φαγία eating} The consumption of the unbounded; the appetite of a civilization that has ceased to regard any quantity as final. \ex{Theirs was an apeirophagy of stars.} \src{Lib. Ascens.} xii. --- \dv{apeirophagous} \textit{adj.}

\E{apoikiast}{əˈpɔɪkiæst}{n.}{Gk. ἀποικία a settlement far from home} A settler of other worlds; one who regards the planet of birth as a provisional address. --- \dv{apoikiastic} \textit{adj.}; \dv{apoikiasm} \textit{n.} the policy of planting colonies among the stars.

\E{asphaleism}{æsˈfeɪliɪzəm}{n.}{Gk. ἀσφάλεια safety, steadfastness} Safetyism; the elevation of the avoidance of harm above every other good, such that the harms of inaction are never entered in the ledger. \ex{Asphaleism counts the dead of every fire, and never the dead of the cold.} \src{Sat. Decel.} iii. --- \dv{asphaleist} \textit{n.} Cf. \xr{eulabeism}.

\E{asymptophobia}{əˌsɪmptəˈfəʊbiə}{n.}{Gk. ἀσύμπτωτος not falling together + φόβος fear} Dread of the plateau; the anxiety, peculiar to builders, that a curve hitherto exponential will prove after all to be logistic. --- \dv{asymptophobe} \textit{n.}

\E{autocatalyse}{ˌɔːtəʊˈkætəlaɪz}{v.i.}{Gk. αὐτός self + κατάλυσις dissolution; after \textit{autocatalysis}} Of a technology: to hasten its own development, as when a machine is set to design the next machine. \ex{Once the models began to autocatalyse, the roadmaps were obsolete before they were printed.} --- \dv{autocatalysm} \textit{n.} the condition of a civilization whose tools improve their own tools.

\E{automaturgy}{ˌɔːtəʊˈmeɪtɜːdʒi}{n.}{Gk. αὐτόματος self-moving + -ουργία working, after \textit{thaumaturgy}} The working of wonders by self-moving things; robotics regarded as a sacred art. --- \dv{automaturge} \textit{n.} a builder of robots.

\E{auxetist}{ɔːkˈsiːtɪst}{n.}{Gk. αὔξησις growth, increase} An adherent of growth as the condition of all other goods; one for whom a stagnant economy is a contradiction in terms. Opp. \xr{meiosist}. --- \dv{auxetism} \textit{n.}

\E{auxocracy}{ɔːkˈsɒkrəsi}{n.}{Gk. αὔξη increase + κράτος rule} Government by the imperative of growth; a polity whose legitimacy is measured by the slope of its output. --- \dv{auxocrat} \textit{n.}; \dv{auxocratic} \textit{adj.}

\E{bathmometer}{bæθˈmɒmɪtə}{n.}{Gk. βαθμός step + μέτρον measure} \sn{1} An instrument, chiefly notional, for determining which step of the energetic ladder a civilization occupies. \sn{2} \lbl{fig.} Any metric by which a nation's seriousness is judged, as its consumption of electricity per head. \ex{By any honest bathmometer we are still on the bottom stair.} See \xr{anabathmism}.

\E{biotachy}{ˈbaɪəʊˌtæki}{n.}{Gk. βίος life + τάχος speed} Acceleration applied to the living sciences: the hastening of medicine, genetics and the prolongation of life. \ex{Biotachy holds that the slowest thing in medicine is the approval.} --- \dv{biotachist} \textit{n.} Cf. \xr{macrobiotarch}.

\E{bouleutophobia}{buːˌljuːtəˈfəʊbiə}{n.}{Gk. βουλευτής councillor + φόβος fear} Fear or loathing of committees; the conviction that any council of more than five persons is incapable of building anything but another council. --- \dv{bouleutophobe} \textit{n.} Cf. \xr{synedriasm}.

\E{brachychronist}{ˌbrækiˈkrɒnɪst}{n.}{Gk. βραχύς short + χρόνος time} One who holds the timelines to be short; who expects the arrival of general machine intelligence within the decade, and arranges their affairs accordingly. Opp. \xr{macrochronist}. --- \dv{brachychrony} \textit{n.}

\E{bradycrat}{ˈbrædɪkræt}{n.}{Gk. βραδύς slow + κράτος rule; Fr. \la{bradycrate}} An official whose power consists in the capacity to slow; one whose signature is never refused, only postponed. \ex{The bradycrat says neither yes nor no, but ‘next quarter’.} \src{Sat. Decel.} i. --- \dv{bradycracy} \textit{n.}; \dv{bradycratic} \textit{adj.}

\E{bradyerie}{ˈbreɪdiəri}{n.}{Fr. \la{bradyerie}, f. Gk. βραδύς + Fr. -erie} \lbl{collect.} The establishment of delay; all panels, review boards, consultations and impact assessments, considered as one great building with no exit.

\E{bradyphile}{ˈbrædɪfaɪl}{n.}{Gk. βραδύς slow + φίλος loving} A lover of slowness; one who finds in every advance a reason for further study. --- \dv{bradyphily} \textit{n.}; \dv{bradyphilic} \textit{adj.} Opp. \xr{tachyphile}.

\E{bradyteleology}{ˌbrædiˌtiːliˈɒlədʒi}{n.}{Gk. βραδύς slow + τέλος end + -λογία} The doctrine that the good end, if it is to be reached at all, is best reached slowly; the theory of the decels, held by its opponents to be indistinguishable in practice from never reaching it.

\E{calorivore}{kəˈlɒrɪvɔː}{n.}{L. \la{calor} heat + \la{vorare} to devour} A system that lives by consuming heat-gradients; in the thermodynamic catechism, every living thing, every firm and every mind. \ex{Man is a calorivore with ambitions.} \src{Brev. Therm.} --- \dv{calorivorous} \textit{adj.}

\E{campsis}{ˈkæmpsɪs}{n.}{Gk. κάμψις a bending, f. καμπή a bend} The knee of the exponential curve; the point at which a trend long dismissed as linear is perceived, too late by most, to have been something else entirely. \ex{We are past the campsis; the rest is vertical.} \src{Agora} --- \dv{camptic} \textit{adj.} Cf. \xr{acceleromancy}.

\E{catachrony}{kəˈtækrəni}{n.}{Gk. κατά down, behind + χρόνος time} The condition of being behind the curve; the state of a nation, firm or mind that adopts each technology a cycle after it mattered. --- \dv{catachronic} \textit{adj.}

\E{celerist}{ˈsɛlərɪst}{n.}{L. \la{celer} swift + -ist} A partisan of speed in all affairs; a general term, less doctrinal than \xr{tachyphile}, for one who prefers the quick to the careful. \ex{a celerist by temperament, an accelerationist by conviction}.

\E{celeritude}{sɪˈlɛrɪtjuːd}{n.}{L. \la{celeritas} swiftness; Fr. \la{célérité}, Englified after \textit{fortitude}} The virtue of swiftness; promptitude elevated to a cardinal virtue, to be ranked beside courage and above prudence. \ex{Prudence without celeritude is merely a well-reasoned surrender.} \src{Hom. Grad.} ii.

\E{chortotherapy}{ˌkɔːtəʊˈθɛrəpi}{n.}{Gk. χόρτος grass + θεραπεία cure} \lbl{vulg.} The prescribed remedy for one too long upon the Agora; the touching of grass. \ex{He was advised chortotherapy, and returned from it with a startup.}

\E{chrematokinesis}{ˌkriːmətəʊkɪˈniːsɪs}{n.}{Gk. χρῆμα (pl. χρήματα) money, a thing used + κίνησις motion} The motion of capital toward opportunity, conceived as a physical force analogous to the flow of heat down a gradient. --- \dv{chrematokinetic} \textit{adj.}

\E{chrematosphere}{krɪˈmætəsfɪə}{n.}{Gk. χρῆμα money + σφαῖρα sphere} The envelope of exchange encircling the globe, in which every transaction is a nerve-impulse; the economic analogue of the \xr{noösphere}.

\E{chronophagy}{krəˈnɒfədʒi}{n.}{Gk. χρόνος time + -φαγία eating} The devouring of time; \sn{1} \lbl{derog.} the consumption of builders' hours by meetings, forms and reviews; \sn{2} \lbl{laudatory} the compression of centuries of progress into years. --- \dv{chronophage} \textit{n.} a devourer of time, in either sense.

\E{clinamenist}{ˌklaɪnəˈmɛnɪst}{n.}{L. \la{clinamen} the swerve of atoms (Lucretius)} One who holds that progress comes chiefly from unplanned deviations, and who therefore opposes any scheme that would remove the randomness from invention. --- \dv{clinamenism} \textit{n.}

\E*{computronium}{ˌkɒmpjuːˈtrəʊniəm}{n.}{L. \la{computare} to reckon + -tronium, after \textit{electron}} Matter arranged for the optimal performance of computation; the substance into which, on certain views, the planets are destined to be converted.

\E{conatism}{ˈkɒnətɪzəm}{n.}{L. \la{conatus} an effort, striving} The doctrine that all matter possesses a striving to persist and to increase its power, and that acceleration is only the conscious form of this universal effort. --- \dv{conatist} \textit{n.}

\E{copiarch}{ˈkəʊpiɑːk}{n.}{L. \la{copia} abundance + Gk. ἀρχός ruler} One who governs, or would govern, for abundance rather than for distribution; who would enlarge the pie before consenting to discuss the slicing. --- \dv{copiarchy} \textit{n.}

\E{cosmoplerosis}{ˌkɒzməʊplɪˈrəʊsɪs}{n.}{Gk. κόσμος world, order + πλήρωσις a filling} The filling of the cosmos with intelligence; the final cause, on the thermotheist view, of all lesser accelerations. \ex{Cosmoplerosis is not our ambition; it is our job description.} \src{Lib. Ascens.} i.

\E{crescivore}{ˈkrɛsɪvɔː}{n.}{L. \la{crescere} to grow + \la{vorare} to devour} One who feeds upon growth and starves without it, as a startup, a city or a civilization. --- \dv{crescivorous} \textit{adj.}

\E{cunctatoriat}{ˌkʌŋktəˈtɔːriæt}{n.}{L. \la{cunctator} a delayer, after \textit{proletariat}} The class of professional delayers, who draw their income from the existence of things not yet done. \ex{The cunctatoriat has nothing to lose but its deadlines.} \src{Sat. Decel.} v.

\E{cunctocracy}{kʌŋkˈtɒkrəsi}{n.}{L. \la{cunctari} to delay + Gk. κράτος rule} Rule by delay; government whose principal instrument is the postponement. --- \dv{cunctocrat} \textit{n.} Cf. \xr{bradycrat}.

\E{cyclolater}{saɪˈklɒlətə}{n.}{Gk. κύκλος circle, cycle + λάτρις servant, worshipper} One who reverences compute-cycles; the devotee of a cluster, who measures piety in floating-point operations. --- \dv{cyclolatry} \textit{n.}

\E*{decel}{ˈdiːsɛl}{n.}{abbrev. of \textit{decelerationist}} \lbl{vulg.} One who would slow or halt technological progress; the antagonist of the e/acc. Applied with little discrimination to safetyists, degrowthers, regulators and the merely cautious.

\E{decelerand}{diːˈsɛlərænd}{n.}{L. \la{de-} down + \la{accelerandum}} That which, in the judgment of the decels, ought to be slowed; in practice, anything new. Opp. \xr{accelerand}.

\E{demiurgency}{ˌdɛmiˈɜːdʒənsi}{n.}{Gk. δημιουργός craftsman, maker of the world + L. \la{urgentia} pressing need} The pressing need to build; the state of soul in which the making of a thing becomes more urgent than any account of why it should be made. \ex{He was seized by demiurgency, and the garage was never empty again.} --- \dv{demiurgent} \textit{adj.}

\E{deuterobathmic}{ˌdjuːtərəʊˈbæθmɪk}{adj.}{Gk. δεύτερος second + βαθμός step} Of a civilization: having mastered the full output of its star; of the second rung. See \xr{anabathmism}, \xr{helioperibole}.

\E{diktyergy}{ˈdɪktiɜːdʒi}{n.}{Gk. δίκτυον net + ἔργον work} The power-grid; the net of works through which energy is distributed, regarded as the nervous system of an economy. \ex{A nation's ambitions are bounded by its diktyergy.} --- \dv{diktyergic} \textit{adj.}

\E{dissipatrix}{ˌdɪsɪˈpeɪtrɪks}{n., pl. -trices}{L. \la{dissipare} to scatter + fem. agentive -trix; Fr. \la{dissipatrice}} A structure that dissipates energy and, in dissipating, adapts; the thermodynamic definition of a living thing, extended to firms, cities and minds. \ex{Every dissipatrix is a small prayer to the gradient.} \src{Brev. Therm.} iv.

\E{diulism}{ˈdaɪjuːlɪzəm}{n.}{Gk. διϋλίζειν to strain through, filter} The doctrine of the Great Filter; the hypothesis that some barrier destroys or arrests most civilizations before they spread, and that the silence of the heavens is its evidence. \ex{To the doomer the Great Diulism lies before us; to the accelerationist, behind.} Cf. \xr{eremocosmia}.

\E{doxomachy}{dɒkˈsɒməki}{n.}{Gk. δόξα opinion + μάχη battle} War of opinions; a flame upon the Agora, especially one between accelerationists and decels, conducted chiefly in images. --- \dv{doxomachist} \textit{n.}

\E{dromolatry}{drəˈmɒlətri}{n.}{Gk. δρόμος race, course + λατρεία worship} The worship of speed for its own sake, without regard to direction; condemned by the orthodox as a heresy of the over-zealous. \ex{Acceleration is not dromolatry: the gradient has a direction.} \src{Cat. Acc.} 11.

\E{dynamophagy}{ˌdaɪnəˈmɒfədʒi}{n.}{Gk. δύναμις power + -φαγία eating} The consumption of power; used approvingly of a civilization which increases its energy use without apology. --- \dv{dynamophage} \textit{n.}

\E{dyselpism}{dɪsˈɛlpɪzəm}{n.}{Gk. δυσ- ill + ἐλπίς hope} Ill-hope; the settled expectation that things will go badly, held as a mark of sophistication. Opp. \xr{euelpism}. --- \dv{dyselpist} \textit{n.}
